% titel: Blatt 5, A5: Vergleich des Atomradius

\subsection*{Frage}
Welches Atom hat je Paar den grösseren Atomradius und warum?

\subsection*{Lösung}
\begin{itemize}
  \item \textbf{P, Cl}: P grösser. Gleiche Periode (3), P steht links von Cl und
    hat die geringere Kernladung $\to$ grösserer Radius.
  \item \textbf{P, As}: As grösser. Gleiche Gruppe (15), As liegt eine Periode
    tiefer $\to$ zusätzliche Schale $\to$ grösserer Radius.
  \item \textbf{P, Al}: Al grösser. Gleiche Periode (3), Al steht links von P
    $\to$ geringere Kernladung $\to$ grösserer Radius (Al: Metall, sehr viel
    grösser als P).
  \item \textbf{P, N}: P grösser. Gleiche Gruppe (15), P liegt unter N $\to$
    eine zusätzliche Schale $\to$ grösserer Radius.
\end{itemize}

\begin{ergebnis}
grösser: P (über Cl), As (über P), Al (über P), P (über N).
\end{ergebnis}

\begin{kontrolle}
Reihenfolge im PSE: zuerst Periode 3 bei P/Al/Cl (links $>$ rechts: Al $>$ P $>$
Cl), dann Gruppe 15 (unten grösser: As $>$ P $>$ N). Beide Regeln liefern
konsistente Grössenverhältnisse.
\end{kontrolle}

\begin{merksatz}
Im Radius-Vergleich: gleiche Periode ist links grösser, gleiche Gruppe ist unten
grösser --- beide Male greift $Z_{\mathrm{eff}}$ bzw.\ die Schalenzahl.
\end{merksatz}
